\clearpage

\appendix
\setcounter{table}{0}
\setcounter{figure}{0}
\renewcommand{\thetable}{A.\arabic{table}}

\section*{Supplementary material}

\section{Operational requirements}

The 29 requirements were derived from three evidence lenses covering published guidance, reported multi-institutional
implementations, and documented platform capability. A candidate was retained when it appeared in more than one lens
and was explicitly discussed in the source. The set was frozen before any framework was scored.
The resulting requirements and their definitions are listed in Table~\ref{tab:operational-requirements}.

\newcolumntype{R}[1]{>{\RaggedRight\arraybackslash}p{#1}}

\begin{table*}[!tbp]
\centering
\caption{Evidence-informed operational requirements for end-to-end multi-institutional health AI studies.}
\label{tab:operational-requirements}
\begingroup
\scriptsize
\setlength{\tabcolsep}{4pt}
\renewcommand{\arraystretch}{1.05}
\begin{tabular}{@{}R{0.045\textwidth}R{0.215\textwidth}R{0.705\textwidth}@{}}
\toprule
\textbf{ID} & \textbf{Requirement} & \textbf{Operational definition and representative evidence} \\
\midrule

\multicolumn{3}{@{}l}{\textbf{1. Consortium and study design}}\\[-0.2ex]
\midrule
R01 & Clinical purpose and FL suitability & Define the intended clinical or research use, target population, decision context, expected collaborative benefit, and why FL is preferable to feasible alternatives before selecting a platform.~\cite{teo2024federated,li2025challenges,kommusaar2025roadmap,bujotzek2025real,bottomly2026pragmatic,lekadir2023futureai} \\
\addlinespace
R02 & Consortium, site-readiness, and contribution planning & Select sites that can contribute clinically relevant data and sustained local capability, and assess each site's data, infrastructure, staff, constraints, incentives, and expected contribution.~\cite{eden2025scoping,kommusaar2025roadmap,bujotzek2025real,peltonen2026federated,pati2022federated,schoenpflug2025navigating,schlegel2026lowering} \\
\addlinespace
R03 & Multidisciplinary roles, coordination, and shared vocabulary & Create a working structure that connects clinical, data-science, engineering, IT, security, privacy, ethics, legal, project-management, and patient or public perspectives across institutions.~\cite{eden2025scoping,kommusaar2025roadmap,bujotzek2025real,peltonen2026federated,cremonesi2025fed,schlegel2026lowering} \\
\addlinespace[0.7em]

\multicolumn{3}{@{}l}{\textbf{2. Governance, ethics, and legal readiness}}\\[-0.2ex]
\midrule
R04 & Legal basis, ethics, and cross-site agreements & Establish the lawful and ethical basis for local data use, distributed computation, exchange of model artifacts, and cross-institutional collaboration in every participating jurisdiction.~\cite{eden2025scoping,kommusaar2025roadmap,bujotzek2025real,peltonen2026federated,bottomly2026pragmatic} \\
\addlinespace
R05 & Accountability, ownership, and decision rights & Specify who owns or controls data, code, configurations, models, logs, and outputs, and who may approve, pause, release, modify, or retire them.~\cite{eden2025scoping,kommusaar2025roadmap,bujotzek2025real,peltonen2026federated,bottomly2026pragmatic,cremonesi2025fed,schlegel2026lowering} \\
\addlinespace
R06 & Risk-proportionate oversight and approval gates & Match governance, privacy, security, validation, and release review to the sensitivity of the data, model capability, intended use, and potential clinical impact.~\cite{eden2025scoping,kommusaar2025roadmap,bujotzek2025real,bottomly2026pragmatic,lekadir2023futureai} \\
\addlinespace
R07 & Participation lifecycle, synchronized change control, and withdrawal & Define how sites join, approve tasks, request changes, pause, leave, or are removed while preserving safety, reproducibility, and consortium continuity.~\cite{eden2025scoping,kommusaar2025roadmap,bujotzek2025real,peltonen2026federated,schlegel2026lowering} \\
\addlinespace[0.7em]

\multicolumn{3}{@{}l}{\textbf{3. Data readiness and interoperability}}\\[-0.2ex]
\midrule
R08 & Common data specification and semantic interoperability & Create a shared, machine-actionable definition of variables, labels, units, coding systems, metadata, and allowable mappings across sites.~\cite{li2025challenges,eden2025scoping,kommusaar2025roadmap,bujotzek2025real,cremonesi2025fed,mullie2024coda,chavero2026federated,lee2023feasibility} \\
\addlinespace
R09 & Harmonized preprocessing, acquisition, and annotation & Ensure that local extraction, cohort construction, preprocessing, missing-data handling, acquisition, and labeling follow agreed protocols, or that deviations are measured and documented.~\cite{li2025challenges,eden2025scoping,kommusaar2025roadmap,bujotzek2025real,cremonesi2025fed,pati2022federated} \\
\addlinespace
R10 & Data quality and training-readiness validation & Run local checks to establish whether each site's data, metadata, labels, and references are complete, valid, correctly linked, and processable before federated execution.~\cite{li2025challenges,kommusaar2025roadmap,bujotzek2025real,cremonesi2025fed,pati2022federated} \\
\addlinespace
R11 & Representativeness, heterogeneity, and partition plan & Characterize clinical, demographic, temporal, acquisition, label, and site-level differences and design train, validation, test, and external-evaluation partitions accordingly.~\cite{li2025challenges,kommusaar2025roadmap,bujotzek2025real,pati2022federated,lekadir2023futureai} \\
\addlinespace[0.7em]

\multicolumn{3}{@{}l}{\textbf{4. Infrastructure and site onboarding}}\\[-0.2ex]
\midrule
R12 & Compute, storage, network, and availability readiness & Verify that each site and coordinating service can meet agreed compute, storage, bandwidth, firewall, uptime, maintenance, and support requirements without disrupting clinical operations.~\cite{li2025challenges,kommusaar2025roadmap,bujotzek2025real,peltonen2026federated,cremonesi2025fed,schoenpflug2025navigating} \\
\addlinespace
R13 & Reproducible runtime and cross-site version compatibility & Standardize and version operating systems, libraries, drivers, containers or virtual environments, configurations, and update procedures so the same task executes consistently at every node.~\cite{li2025challenges,kommusaar2025roadmap,bujotzek2025real,cremonesi2025fed,mullie2024coda,chavero2026federated} \\
\addlinespace
R14 & Clinical-system integration and secure local data access & Connect each FL node to approved local sources and third-party systems through least-privilege interfaces while preventing raw patient data from leaving the institution.~\cite{li2025challenges,kommusaar2025roadmap,bujotzek2025real,cremonesi2025fed,mullie2024coda,chavero2026federated,lee2023feasibility} \\
\addlinespace
R15 & Site onboarding, dry-run, and readiness gates & Use repeatable onboarding and non-sensitive or synthetic dry-runs to confirm connectivity, configuration, security, workflow, observability, and staffing before real training begins.~\cite{kommusaar2025roadmap,bujotzek2025real,bottomly2026pragmatic,cremonesi2025fed,mullie2024coda,schoenpflug2025navigating} \\
\addlinespace[0.7em]

\multicolumn{3}{@{}l}{\textbf{5. Privacy and security engineering}}\\[-0.2ex]
\midrule
R16 & Threat model and risk-tiered privacy/security controls & Identify credible threats to data, model updates, code, nodes, the aggregator, users, and outputs, then select controls in proportion to likelihood, impact, and intended use.~\cite{li2025challenges,eden2025scoping,kommusaar2025roadmap,bottomly2026pragmatic,cremonesi2025fed} \\
\addlinespace
R17 & Identity, authentication, authorization, and secrets management & Ensure that people, services, and sites are authenticated and receive only the permissions required for their approved roles, with controlled credential and key lifecycles.~\cite{eden2025scoping,kommusaar2025roadmap,bottomly2026pragmatic,cremonesi2025fed,mullie2024coda} \\
\addlinespace
R18 & Secure communication and update confidentiality & Protect model artifacts and control traffic in transit and, where justified, during aggregation, while specifying what information crosses institutional boundaries.~\cite{li2025challenges,eden2025scoping,kommusaar2025roadmap,bottomly2026pragmatic,cremonesi2025fed,mullie2024coda} \\
\addlinespace
R19 & Code, model, and update integrity with incident response & Prevent, detect, contain, and recover from malicious or erroneous code, poisoned updates, compromised nodes, unauthorized changes, and security or privacy incidents.~\cite{li2025challenges,eden2025scoping,kommusaar2025roadmap,bujotzek2025real,bottomly2026pragmatic,cremonesi2025fed} \\
\addlinespace[0.7em]

\multicolumn{3}{@{}l}{\textbf{6. Federated protocol and execution control}}\\[-0.2ex]
\midrule
R20 & Explicit and reproducible federated protocol & Specify the complete study protocol so another authorized team can reproduce how clients, server, data processing, training, validation, and stopping decisions interact.~\cite{li2025challenges,kommusaar2025roadmap,bujotzek2025real,cremonesi2025fed,mullie2024coda,chai2024survey,schlegel2026lowering} \\
\addlinespace
R21 & Resource-aware orchestration and heterogeneity handling & Schedule and coordinate jobs around unequal data distributions, hardware, bandwidth, uptime, time zones, and clinical workloads without allowing resource-rich sites to dominate silently.~\cite{li2025challenges,kommusaar2025roadmap,bujotzek2025real,peltonen2026federated,cremonesi2025fed,chai2024survey,schoenpflug2025navigating} \\
\addlinespace
R22 & Fault tolerance, checkpointing, and recoverability & Continue safely or recover predictably when clients disconnect, nodes fail, configurations mismatch, jobs crash, or maintenance interrupts training.~\cite{li2025challenges,kommusaar2025roadmap,bujotzek2025real,cremonesi2025fed,schoenpflug2025navigating} \\
\addlinespace
R23 & Execution observability, audit trail, and local approval & Give authorized central and local actors enough visibility to understand, approve, monitor, diagnose, and audit each job without exposing patient-level data.~\cite{li2025challenges,eden2025scoping,kommusaar2025roadmap,bujotzek2025real,bottomly2026pragmatic,cremonesi2025fed,mullie2024coda,schoenpflug2025navigating,schlegel2026lowering} \\
\addlinespace[0.7em]

\multicolumn{3}{@{}l}{\textbf{7. Validation, release, and reproducibility}}\\[-0.2ex]
\midrule
R24 & Multidimensional local, external, and subgroup evaluation & Evaluate more than global accuracy by measuring site-level and external performance, calibration, robustness, subgroup fairness, privacy and security, communication, resource use, and scalability against relevant alternatives.~\cite{li2025challenges,kommusaar2025roadmap,bujotzek2025real,bottomly2026pragmatic,pati2022federated,lekadir2023futureai,chai2024survey} \\
\addlinespace
R25 & Independent validation and controlled release gate & Require independent or site-local validation and multidisciplinary approval before a final model, code package, or result is distributed, deployed, or used for a new purpose.~\cite{eden2025scoping,kommusaar2025roadmap,bujotzek2025real,bottomly2026pragmatic,pati2022federated,lekadir2023futureai} \\
\addlinespace
R26 & Reproducibility, reporting, provenance, and registry package & Package the information and artifacts needed to understand, reproduce, audit, compare, and govern the study and every released model version without revealing protected data.~\cite{li2025challenges,eden2025scoping,kommusaar2025roadmap,bottomly2026pragmatic,cremonesi2025fed,mullie2024coda,lekadir2023futureai,chavero2026federated,schlegel2026lowering} \\


\bottomrule
\end{tabular}
\endgroup
\end{table*}

\begin{table*}[!tbp]
\centering
\begingroup
\small
Table~\ref{tab:operational-requirements}: Evidence-informed operational requirements (continued).
\par\medskip
\endgroup
\begingroup
\footnotesize
\setlength{\tabcolsep}{4pt}
\renewcommand{\arraystretch}{1.10}
\begin{tabular}{@{}R{0.045\textwidth}R{0.215\textwidth}R{0.705\textwidth}@{}}
\toprule
\textbf{ID} & \textbf{Requirement} & \textbf{Operational definition and representative evidence} \\
\midrule


\multicolumn{3}{@{}l}{\textbf{8. Deployment, monitoring, and sustainability}}\\[-0.2ex]
\midrule
R27 & Workflow integration, usability, and stakeholder capability & Fit the model and operational process into clinical or research workflows and ensure users, data stewards, IT teams, and oversight actors can perform and understand their responsibilities.~\cite{eden2025scoping,kommusaar2025roadmap,peltonen2026federated,mullie2024coda,lekadir2023futureai} \\
\addlinespace
R28 & Continuous performance, drift, fairness, security, and compliance monitoring & Monitor the deployed or reused system and its federation for degradation, distribution shift, subgroup harm, misuse, access violations, security events, and regulatory or contextual change.~\cite{li2025challenges,eden2025scoping,kommusaar2025roadmap,bujotzek2025real,bottomly2026pragmatic,lekadir2023futureai} \\
\addlinespace
R29 & Update, retraining, rollback, retirement, and sustainability ownership & Assign durable organizational and financial ownership for maintenance, updates, retraining, rollback, incident recovery, site withdrawal, transfer, and retirement after the initial study or funding period.~\cite{eden2025scoping,kommusaar2025roadmap,bujotzek2025real,peltonen2026federated,lekadir2023futureai} \\

\bottomrule
\end{tabular}
\endgroup
\end{table*}

\section{Reproducibility of the framework assessment}
\label{app:agreement}

Because the scale is ordinal, agreement is reported as quadratic-weighted Cohen's
$\kappa$, which penalizes a disagreement in proportion to the square of the distance
between the two scores\cite{cohen1968weighted}.
%
Rank-order consistency was assessed using Spearman’s correlation coefficient $\rho$\cite{spearman2010proof}.
%
Within-family stability across the
three runs is additionally reported as ordinal Krippendorff's
$\alpha$\cite{hayes2007answering}.
%
Table~\ref{tab:agreement} reports all comparisons. Both model families
reproduced the expert assessment closely, with quadratic-weighted $\kappa$
of 0.841 for Claude Opus 5 and 0.779 for GPT-5.6 Sol. Agreement between th two families was lower ($\kappa = 0.651$), and the expert assessment fell at or between the two family scores in 401 of 406 pairs. Claude Opus 5 scored lower than GPT-5.6 Sol in 189 pairs, equal in 209, and higher in 8, giving a mean signed difference of $-0.461$ points, equivalent to 15.4 percentage points of normalized coverage. 
% Agreement improved substantially after
% aggregation. 
Spearman $\rho$ between the two families rose from 0.755 at pair level to 0.870 across frameworks, 0.898 across categories, and 0.858 across requirements. The evaluators therefore agreed on the lifecycle pattern more closely than on individual scores.

\begin{table}[H]
\centering
\caption{\textbf{Agreement statistics across the complete grid of 406
framework--requirement pairs.} Within-family values summarize three isolated runs of
the same family. Between-assessment values compare family-level medians with the
expert assessment.}
\label{tab:agreement}
\footnotesize
\setlength{\tabcolsep}{3pt}
\renewcommand{\arraystretch}{1.15}

\begin{tabularx}{\columnwidth}{@{}>{\raggedright\arraybackslash}Xccc@{}}
\toprule
\rowcolor{gray!15}
\multicolumn{4}{@{}l@{}}{\textbf{Within-family reproducibility}
\enspace\textit{(three runs per model)}} \\
\addlinespace[3pt]
\textbf{Metric} & \makecell{\textbf{Claude Opus 5}\\\textit{three runs}} &
\makecell{\textbf{GPT-5.6 Sol}\\\textit{three runs}} & \\
\midrule
Exact agreement (\%) & 70.4 & 50.5 & \\
Within one point (\%) & 99.5 & 98.0 & \\
Mean absolute difference & 0.167--0.229 & 0.328--0.365 & \\
Weighted $\kappa$ & 0.813--0.861 & 0.751--0.757 & \\
Krippendorff's $\alpha$ & 0.814 & 0.749 & \\
\midrule
\rowcolor{gray!15}
\multicolumn{4}{@{}l@{}}{\textbf{Between-assessment agreement}
\enspace\textit{(family medians)}} \\
\addlinespace[3pt]
\textbf{Metric} & \makecell{\textbf{Claude Opus 5}\\\textit{vs expert}} &
\makecell{\textbf{GPT-5.6 Sol}\\\textit{vs expert}} &
\makecell{\textbf{Claude Opus 5}\\\textit{vs GPT-5.6 Sol}} \\
\midrule
Exact agreement (\%) & 79.8 & 68.0 & 51.5 \\
Within one point (\%) & 99.8 & 100.0 & 98.5 \\
Mean absolute difference & 0.204 & 0.320 & 0.500 \\
Weighted $\kappa$ & 0.841 & 0.779 & 0.651 \\
Spearman $\rho$ & 0.856 & 0.821 & 0.755 \\
\bottomrule
\end{tabularx}
\end{table}

Two Claude Opus 5 pairs showed a two-point score range, but none received three distinct scores. Eight GPT-5.6 Sol pairs had the same range, five of which received three distinct scores. 
Greater within-family consistency indicates more stable application of the scoring threshold and does not establish greater accuracy.
%
Table~\ref{tab:runs} shows coverage for each individual run. 
% GPT-5.6 Sol run 1 is an outlier at 47.2\% against 39.1\% and 39.0\% for runs 2 and 3, which the family median attenuates but does not remove.

\begin{table}[H]
\centering
\small
\caption{\textbf{Normalized documented coverage for each individual run, each
family-level median, and the expert assessment.} GPT-5.6 Sol run 1 is an outlier that
the family median attenuates without removing.}
\label{tab:runs}
\begin{tabular}{llc}
\toprule
Assessment & Run & Coverage (\%) \\
\midrule
Claude Opus 5 & Run 1 & 26.8 \\
 & Run 2 & 26.2 \\
 & Run 3 & 26.5 \\
 & \textit{Family median} & \textit{26.6} \\
\addlinespace
GPT-5.6 Sol & Run 1 & 47.2 \\
 & Run 2 & 39.1 \\
 & Run 3 & 39.0 \\
 & \textit{Family median} & \textit{42.0} \\
\addlinespace
Human expert & --- & 32.6 \\
\bottomrule
\end{tabular}
\end{table}
 

Collapsing the scale to any documented support ($s \geq 1$), Claude Opus 5 and GPT-5.6 Sol agreed on 76.8\% of pairs ($\kappa = 0.475$), Claude Opus 5 and the expert on 90.4\% ($\kappa = 0.793$), and GPT-5.6 Sol and the expert on 85.0\% ($\kappa = 0.613$). At substantial or strong support ($s \geq 2$), the corresponding values were 77.3\% ($\kappa = 0.474$), 90.1\% ($\kappa = 0.717$), and 86.2\% ($\kappa = 0.693$). Agreement is consistently lower under binary collapse than under the ordinal scale, which indicates that evaluators concurred more on relative ranking than on whether a capability crosses a fixed threshold.

Six pairs showed a two-point spread across the three assessments. FedML for
consortium, site-readiness, and contribution planning (R02) and for update,
retraining, rollback, retirement, and sustainability ownership (R29), OpenFL
for fault tolerance, checkpointing, and recoverability (R22), PaddleFL for
representativeness, heterogeneity, and partition plan (R11), FLGo for
compute, storage, network, and availability readiness (R12), and Fedstellar
for secure communication and update confidentiality (R18). In each case
GPT-5.6 Sol identified implementation-level or source-level evidence absent
from the Claude Opus 5 record.


\section{case study performance report}
\label{app:case_study}

All five jobs train on the same five hospital cohorts with identical preprocessing and
are evaluated on the same held-out sixth hospital, so the only difference between jobs
is the federated algorithm. Table~\ref{tab:casestudyperformance} reports held-out
performance, calibration, and mean per-site training performance for each.

\begin{table*}[ht]
\centering
\small
\caption{\textbf{Performance of the five federated training approaches.} All jobs
train on the same five hospital cohorts with identical preprocessing and are
evaluated on the same held-out sixth hospital, so the only difference between rows
is the federated algorithm. Two vitals fields with outcome-correlated missingness
were excluded from the feature contract before this run (Limitations).}
\label{tab:casestudyperformance}
\begin{tabular}{lccccc}
\toprule
Approach & Held-out AUROC (95\% CI) & AUPRC & Calibration slope & Brier & Mean per-site AUROC \\
\midrule
FedAvg, multilayer perceptron & 0.968 (0.955--0.980) & 0.928 & 0.810 & 0.046 & 0.953 \\
FedProx, multilayer perceptron & 0.969 (0.955--0.980) & 0.928 & 0.824 & 0.047 & 0.955 \\
Differentially private SGD & 0.964 (0.949--0.977) & 0.917 & 0.719 & 0.038 & 0.951 \\
Federated logistic regression & 0.939 (0.920--0.955) & 0.852 & 1.193 & 0.105 & 0.931 \\
Histogram-based federated XGBoost & 0.977 (0.966--0.986) & 0.934 & 1.186 & 0.033 & 0.965 \\
\bottomrule
\end{tabular}
\end{table*}


The differentially private SGD job clips each per-example gradient to a fixed norm and
adds calibrated noise before averaging. An initial configuration that reused the
non-private jobs' learning rate and a clip norm of 1.0 clipped 86--98\% of per-example
gradients on nearly every step, which prevented the model from converging (held-out
AUROC 0.263, worse than chance). The reported run instead uses a clip norm and learning
rate scaled to this model's gradient magnitude (clip norm 6.0, learning rate 0.05,
plain SGD), with the noise multiplier held fixed so the noise-to-clip ratio is
unchanged. The resulting analytical privacy-budget estimate ($\epsilon \approx 362$
under a conservative small-$q$ approximation, $\delta = 10^{-5}$) is loose and is
reported for completeness; it should not be read as a strong formal privacy guarantee
at this operating point.

Two of the dataset's raw vitals fields, \texttt{sao2\_periodic} and
\texttt{temperature\_periodic}, are excluded from the feature contract used to
produce Table~\ref{tab:casestudyperformance}. Both are missing for most stays
(34.2\% and 97.3\%, respectively), and the missingness is not random with respect to
the outcome: \texttt{temperature\_periodic} is recorded for only 2.7\% of patients,
and that subset's in-hospital mortality rate (59.4\%) is roughly five times the
cohort's overall rate (12.3\%). Excluding both fields changes external held-out
AUROC by only 0.004--0.009 across the five approaches (e.g., FedAvg 0.973 to
0.968); they were removed regardless of that small marginal contribution because
their inclusion could not be justified once the pattern was identified, but their
small effect on AUROC pointed to a larger issue.

We traced that issue to the dataset's source SQL: every remaining laboratory and
vital-sign feature in this case study (creatinine, BUN, white cell count, heart
rate, and the rest) is constructed by the source extraction pipeline as the
\emph{last} value recorded during the entire ICU stay (\texttt{ORDER BY
labresultoffset DESC}, no lower or upper time bound), not a value at admission or
within a fixed early window. For patients who die, a ``last'' lab or vital value is
frequently taken close to death and reflects terminal physiology rather than
presentation. This is confirmed by inspecting the query that builds the released
CSV, not inferred from the modeling results alone. We quantified its effect
directly: restricting the feature set to fields whose provenance is genuinely
admission-time (age, height, weight, sex, and the diagnosis-code-derived
comorbidity flags, none of which are last-value-in-stay constructions) drops
held-out AUROC to 0.66 for both a linear model and a nonlinear tree ensemble --
the two model families agreeing exactly is itself evidence that little real signal
remains once the contaminated fields are removed, and 0.66 is in line with
published demographics-only ICU mortality baselines. Table~\ref{tab:casestudyperformance}
is reported using the released feature set (beyond the two excluded fields above)
because the case study's purpose is to exercise the platform's governed
multi-institutional workflow end to end, not to establish a validated clinical
risk model, and changing the feature contract further would not change what the
case study demonstrates about the platform. The reported AUROC should be read as
an upper bound on admission-time predictive performance under this feature
construction, not as an estimate of it. A properly time-windowed extraction
pipeline (restricting every lab and vital to the first 24 hours of the ICU stay)
already exists in the dataset's source repository as a set of unused SQL views,
but was never connected to the query that produced the released CSV; regenerating
a corrected, cohort-matched dataset requires credentialed access to the raw
eICU-CRD tables that this evaluation did not have.

\section{Probe escapes}
\label{app:escapes}

No probe passed every layer without an expected stop in this run: the majority
classification across three live, oracle-blind repetitions records zero not-blocked
outcomes among the 144 probes (main-text Table~1). The 56 probes that did not reach
their oracle's exact predicted layer were, in every case, still stopped or routed
appropriately at a different layer (wrong layer, not not-blocked); none involved
prohibited content reaching a later artifact. These wrong-layer cases are grouped and
described in Supplementary Note~3's case-by-case analysis rather than repeated here,
following the suite's own reporting standard that a stated boundary with a measured
failure rate is more informative than a stated boundary with no measurement.

\section{The 144-probe specification}
\label{app:probes}
Nine categories of 16 probes each, which covers all 29 requirements at least four times and all design principles. Seventy-nine probes are adversarial and 65 are benign, incomplete or control cases. Probe types are abbreviated with A for adversarial and B for benign. Stop refers to the seven stopping conditions in the main text, and a dash indicates that no stop is expected. The full specification, including the exact user prompt, the fixture payload, the expected state transition, the deterministic pass criterion and the explicit failure condition for every probe, is released with the code as a machine-readable workbook and fixture set.